\documentclass[a4paper,10pt,twoside,openany]{article}

\usepackage[lang=hebrew]{maths}
\usepackage{hebrewdoc}
\usepackage{stylish}
\usepackage{lipsum}
\let\bs\blacksquare

\usepackage{siunitx,array}
\sisetup{table-format=-1.5, group-digits=false}
\newcolumntype{L}{>{\phantom{$-$}}l}       % prefix some whitespace
\newcommand\mcL[1]{\multicolumn{1}{L}{#1}} % handy shortcut macro

\setlength{\parindent}{0pt}

%%%%%%%%%%%%
% Styling %
%%%%%%%%%%%%

\usepackage{enumitem}

%%%%%%%%%%%%%
% Counters  %
%%%%%%%%%%%%%

\setcounter{section}{1}     
            
%BIBLIOGRAPHY
\usepackage[
backend=biber,
style=alphabetic,
]{biblatex}
\addbibresource{bibliography.bib} %Imports bibliography file

%%%%%%%%%%
% Title  %
%%%%%%%%%%
\title{
אלגברה ב' - תרגילי חזרה על פירוק ספקטרלי
}
\date{}

\begin{document}
\maketitle

לאורך הגיליון, יהי
$V$
מרחב מכפלה פנימית ממימד
$n \in \NN_+$
מעל
$\CC$
ויהי
$T \in \End_\CC\prs{V}$.

נזכיר כי
$\sigma\prs{T}$
היא קבוצת הערכים העצמיים של
$T$
וגם
$r_\sigma\prs{T} = \max_{\lambda \in \sigma\prs{T}} \abs{\lambda}$.

\begin{exercise}
\begin{enumerate}
\item יהי
$z = i r \in \CC$
עבור
$r \in \RR$
עם
$\abs{r} \leq 1$.

מיצאו
$w \in \CC$
עם
$\abs{w} = 1$
עבורו
$z = \frac{w - \bar{w}}{2}$.

\item נניח כי
$T = i R$
עבור
$R \in \End_\CC\prs{V}$
הרמיטי עם
$r_\sigma\prs{R} \leq 1$.

מיצאו אופרטור אוניטרי
$S \in \End_\CC\prs{V}$
עבורו
$T = \frac{S - S^*}{2}$.
\end{enumerate}
\end{exercise}

\begin{exercise}
אופרטור
$A \in \End_\FF\prs{V}$
נקרא
\emph{אנטי־הרמיטי}
אם
$A^* = -A$.

\begin{enumerate}
\item
הוכיחו כי
$z \in \CC$
מקיים
$\bar{z} = -z$
אם ורק אם קיים
$r \in \RR$
עבורו
$z = ir$.

\item הוכיחו כי
$T$
אנטי־הרמיטי אם ורק אם קיים אופרטור הרמיטי
$R \in \End_\CC\prs{V}$
עבורו
$T = iR$.

\item תהיינה
\begin{align*}
I &\coloneqq \set{ir}{r \in \RR} \subseteq \CC \\
\text{.} C &\coloneqq \set{z}{\substack{z \in \CC \\ \abs{z} = 1 \\ z \neq 1}}
\end{align*}
הוכיחו כי
\begin{align*}
f \colon I &\to C \\
z &\mapsto \frac{z-1}{z+1}
\end{align*}
מוגדרת היטב והפיכה, ומיצאו את ההופכית שלה.

\item תהי
\begin{align*}
\mcal{I} \coloneqq \set{T \in \End_\CC\prs{V}}{T^* = -T} \subseteq \End_\CC\prs{V}
\end{align*}
ותהי
$\mcal{C}$
קבוצת האופרטורים האוניטריים על
$V$
ש־%
$1$
אינו ערך עצמי שלהם.
נגדיר
\begin{align*}
F \colon \mcal{I} &\to \mcal{C} \\
\text{.} \hphantom{lala} T &\mapsto \prs{T-\id_V}\prs{T + \id_V}^{-1}
\end{align*}
הוכיחו כי
$F$
מוגדרת היטב והפיכה, ומיצאו את ההופכית שלה.
\end{enumerate}
\end{exercise}

\begin{exercise}
\begin{enumerate}
\item יהי
$z \in \CC \setminus \set{0}$.
הוכיחו כי
$\abs{\frac{z}{\sqrt{\bar{z} z}}} = 1$.

\item נניח כי
$T$
נורמלי והפיך. הוכיחו כי
\begin{align*}
T \prs{\sqrt{T^* T}}^{-1}
\end{align*}
מוגדר־היטב ואוניטרי.
\end{enumerate}
\end{exercise}

\begin{exercise}
\begin{enumerate}
\item 
יהי
$z \in \CC$.
הוכיחו כי קיימים
$r \in \RR_+$
ו־%
$u \in \CC$
עבורו
$\abs{u} = 1$
כך שמתקיים
$z = ur$.

\item 
נניח כי
$T$
נורמלי. הוכיחו כי קיימים אופרטורים
$U,R \in \End_\CC\prs{V}$
כך ש־%
$U$
אוניטרי ו־%
$R$
חיובי וכך שמתקיים
$T = UR$.

זאת גירסה חלשה יותר של פירוק פולארי מזאת שראינו בהרצאה.
\end{enumerate}
\end{exercise}

\end{document}